\documentclass[11pt, a4paper]{article}

\usepackage[T1]{fontenc}
\usepackage[utf8]{inputenc}
\usepackage{tgpagella}
\usepackage[spanish, es-noshorthands]{babel}
\usepackage{amsmath, amssymb, bm}
\usepackage{tabularx}
\usepackage{booktabs}
\usepackage[most]{tcolorbox}
\usepackage[margin=2.5cm]{geometry}
\usepackage{ragged2e}
\usepackage{fancyhdr}

\pagestyle{fancy}
\fancyhf{}
\setlength{\headheight}{16pt}
\lhead{\textbf{UCB Sede Tarija} | Ing. Mecatrónica}
\chead{}
\rhead{IMT-121 Circuitos Electrónicos I | Reconciliación Hito 2}
\lfoot{Prof: M.Sc. Ing. Bernardo Quiroga Turdera}
\rfoot{Página \thepage}
\renewcommand{\headrulewidth}{0.4pt}

\definecolor{ucbblue}{RGB}{0, 51, 102}

\begin{document}

\begin{tcolorbox}[colback=ucbblue!5!white, colframe=ucbblue, title=\textbf{Matriz de Reconciliación de Evidencias (Hito 2)}]
    \centering \textbf{Uso Exclusivo del Instructor}
\end{tcolorbox}

\noindent \textit{Ponderación oficial: Desempeño Práctico (20\%) + Informe IEEE A+B (30\%) + Examen Individual (50\%). La consistencia determina si se activa la microdefensa.}

\vspace{0.3cm}
\renewcommand{\arraystretch}{2}
\noindent
\begin{tabularx}{\textwidth}{
  >{\RaggedRight\arraybackslash\hsize=1.8\hsize}X 
  >{\centering\arraybackslash\hsize=0.6\hsize}X 
  >{\centering\arraybackslash\hsize=0.7\hsize}X 
  >{\centering\arraybackslash\hsize=0.7\hsize}X 
  >{\centering\arraybackslash\hsize=0.7\hsize}X 
  >{\centering\arraybackslash\hsize=0.8\hsize}X 
  >{\centering\arraybackslash\hsize=1.2\hsize}X 
  >{\RaggedRight\arraybackslash\hsize=1.5\hsize}X}
\toprule
\textbf{Estudiante} & \textbf{P(20\%)} & \textbf{R(30\%)} & \textbf{E(50\%)} & \textbf{H2 Pnd.} & \textbf{Consist.} & \textbf{Acción} & \textbf{Nota de Evidencia} \\ 
\midrule
 & & & & & & & \\ \midrule
 & & & & & & & \\ \midrule
 & & & & & & & \\ \midrule
 & & & & & & & \\ \midrule
 & & & & & & & \\ \midrule
 & & & & & & & \\ \midrule
 & & & & & & & \\ 
\bottomrule
\end{tabularx}

\vspace{0.5cm}
\textbf{Claves de Consistencia y Acción:}
\begin{itemize}
    \item \textbf{CONSISTENT:} Evidencias alineadas lógicamente. $\rightarrow$ \textit{Close EC2}.
    \item \textbf{REVIEW:} Diferencias menores o evidencia omitida. $\rightarrow$ \textit{Revisar sin Microdefensa}.
    \item \textbf{DISCREPANT:} Conflicto agudo (ej. Examen deficiente vs. Reporte perfecto). $\rightarrow$ \textit{Microdefensa Selectiva}.
\end{itemize}

\end{document}
